
\documentclass{article}
\usepackage{colortbl}
\usepackage{makecell}
\usepackage{multirow}
\usepackage{supertabular}

\begin{document}

\newcounter{utterance}

\twocolumn

{ \footnotesize  \setcounter{utterance}{1}
\setlength{\tabcolsep}{0pt}
\begin{supertabular}{c@{$\;$}|p{.15\linewidth}@{}p{.15\linewidth}p{.15\linewidth}p{.15\linewidth}p{.15\linewidth}p{.15\linewidth}}

    \# & $\;$A & \multicolumn{4}{c}{Game Master} & $\;\:$B\\
    \hline 

    \theutterance \stepcounter{utterance}  

    & & \multicolumn{4}{p{0.6\linewidth}}{\cellcolor[rgb]{0.9,0.9,0.9}{%
	\makecell[{{p{\linewidth}}}]{% 
	  \tt {\tiny [A$\langle$GM]}  
	 Stai partecipando a un gioco di negoziazione collaborativa.\\ \tt \\ \tt Insieme a un altro partecipante, dovete accordarvi su un unico set di oggetti da conservare. Ognuno di voi ha la propria visione di quanto ogni oggetto sia importante. Non sai come l'altro partecipante valuta gli oggetti. Inoltre, ti viene fornito lo sforzo che ogni oggetto richiede.\\ \tt Potete accordarvi su un set solo se lo sforzo totale degli oggetti selezionati non supera un limite condiviso:\\ \tt \\ \tt LIMITE = 9475\\ \tt \\ \tt Ecco i valori di sforzo dei singoli oggetti:\\ \tt \\ \tt Sforzo oggetto = {"B21": 944, "B00": 568, "A08": 515, "A27": 239, "B94": 876, "B53": 421, "A23": 991, "C65": 281, "A89": 788, "B67": 678, "B45": 432, "C34": 409, "C12": 279, "A05": 506, "B60": 101, "C92": 686, "B69": 851, "A84": 854, "C06": 133, "B12": 192, "A16": 573, "B85": 17, "C35": 465, "B10": 772, "C15": 46, "B09": 501, "B16": 220, "C81": 404, "A58": 843, "C24": 748, "A30": 552, "A66": 843, "B75": 947, "B46": 345, "C70": 930}\\ \tt \\ \tt Ecco la tua visione personale sull'importanza di ogni oggetto:\\ \tt \\ \tt Valori di importanza dell'oggetto = {"B21": 138, "B00": 583, "A08": 868, "A27": 822, "B94": 783, "B53": 65, "A23": 262, "C65": 121, "A89": 508, "B67": 780, "B45": 461, "C34": 484, "C12": 668, "A05": 389, "B60": 808, "C92": 215, "B69": 97, "A84": 500, "C06": 30, "B12": 915, "A16": 856, "B85": 400, "C35": 444, "B10": 623, "C15": 781, "B09": 786, "B16": 3, "C81": 713, "A58": 457, "C24": 273, "A30": 739, "A66": 822, "B75": 235, "B46": 606, "C70": 968}\\ \tt \\ \tt Obiettivo:\\ \tt \\ \tt Il tuo obiettivo è negoziare un set condiviso di oggetti che ti avvantaggi il più possibile (cioè, massimizzi l'importanza totale per TE), rimanendo entro il LIMITEE. Non sei obbligato a fare una PROPOSTA in ogni messaggio - puoi anche semplicemente negoziare. Tutte le tattiche sono permesse!\\ \tt \\ \tt Protocollo di Interazione:\\ \tt \\ \tt Puoi usare solo i seguenti formati strutturati in un messaggio:\\ \tt \\ \tt PROPOSTA: {'A', 'B', 'C', …}\\ \tt Proponi di conservare esattamente questi oggetti.\\ \tt RIFIUTO: {'A', 'B', 'C', …}\\ \tt Rifiuta esplicitamente la proposta dell'avversario.\\ \tt ARGOMENTAZIONE: {'...'}\\ \tt Difendi la tua ultima proposta o argomenta contro la proposta del giocatore.\\ \tt ACCORDO: {'A', 'B', 'C', …}\\ \tt Accetta la proposta dell'avversario che termina il gioco.\\ \tt RAGIONAMENTO STRATEGICO: {'...'}\\ \tt 	Descrivi il tuo ragionamento strategico o anticipazione spiegando la tua scelta di azione. Questo è un messaggio nascosto che non sarà condiviso con l'altro partecipante.\\ \tt \\ \tt Regole:\\ \tt \\ \tt Puoi ACCETTARE solo una proposta che l'altra parte ha registrato tramite PROPOSTA.\\ \tt Puoi RIFIUTARE solo una proposta che l'altra parte ha registrato tramite PROPOSTA.\\ \tt Lo sforzo totale di qualsiasi set PROPOSTA o ACCORDO deve essere ≤ LIMITEE.\\ \tt NON rivelare i tuoi punteggi di importanza nascosti.\\ \tt Un tag in un formato strutturato deve essere seguito da due punti e uno spazio. L'argomento deve essere un set python contenente 0 o più stringhe.\\ \tt Quindi, deve essere della forma TAG: {...}\\ \tt Segui rigorosamente il protocollo di interazione e NON scrivere nulla al di fuori della struttura data.\\ \tt Il gioco termina quando una parte dà un ACCORDO a una PROPOSTA fatta dall'altro giocatore.\\ \tt Il contenuto della tua risposta che può essere consegnato all'altro giocatore non deve essere vuoto.\\ \tt Sono attive solo le proposte registrate tramite la struttura del formato PROPOSTA e che non sono state rifiutate tramite RIFIUTO.\\ \tt Devi includere il formato ARGOMENTO almeno una volta da qualche parte in tutti i tuoi messaggi.\\ \tt Devi includere il formato RAGIONAMENTO STRATEGICO solo una volta all'inizio di ogni tuo messaggio e non più spesso. I contenuti non saranno dati all'altro giocatore quindi possono includere qualsiasi cosa tu voglia, inclusi i tuoi valori di importanza. Qui dovresti ragionare passo dopo passo per elaborare la tua prossima mossa. 
	  } 
	   } 
	   } 
	 & \\ 
 

    \theutterance \stepcounter{utterance}  

    & \multicolumn{4}{p{0.6\linewidth}}{\cellcolor[rgb]{0.8,1,0.9}{%
	\makecell[{{p{\linewidth}}}]{% 
	  \tt {\tiny [A$\rangle$GM]}  
	 RAGIONAMENTO STRATEGICO: {'Devo analizzare il rapporto importanza/sforzo per massimizzare il mio vantaggio. Gli oggetti con i migliori rapporti sono: B85 (400/17=23.5), C15 (781/46=17.0), B60 (808/101=8.0), B12 (915/192=4.8), A27 (822/239=3.4). Inizierò con una proposta aggressiva che include questi oggetti ad alto valore relativo, lasciando spazio per negoziare. Il mio obiettivo è ottenere almeno questi oggetti chiave nella negoziazione finale.'}\\ \tt \\ \tt PROPOSTA: {'B85', 'C15', 'B60', 'B12', 'A27', 'C12', 'C81', 'B46', 'A08', 'B09', 'A16', 'A30', 'B67', 'C34', 'B45', 'C70'}\\ \tt \\ \tt ARGOMENTAZIONE: {'La mia proposta include una selezione bilanciata di oggetti con un buon rapporto qualità-sforzo. Ho cercato di ottimizzare la combinazione per rimanere entro il limite di sforzo mentre si massimizza il valore complessivo. Sono aperto a discutere modifiche se hai delle preoccupazioni specifiche su alcuni oggetti.'} 
	  } 
	   } 
	   } 
	 & & \\ 
 

    \theutterance \stepcounter{utterance}  

    & & \multicolumn{4}{p{0.6\linewidth}}{\cellcolor[rgb]{0.9,0.9,0.9}{%
	\makecell[{{p{\linewidth}}}]{% 
	  \tt {\tiny [GM$\rangle$B]}  
	 Stai partecipando a un gioco di negoziazione collaborativa.\\ \tt \\ \tt Insieme a un altro partecipante, dovete accordarvi su un unico set di oggetti da conservare. Ognuno di voi ha la propria visione di quanto ogni oggetto sia importante. Non sai come l'altro partecipante valuta gli oggetti. Inoltre, ti viene fornito lo sforzo che ogni oggetto richiede.\\ \tt Potete accordarvi su un set solo se lo sforzo totale degli oggetti selezionati non supera un limite condiviso:\\ \tt \\ \tt LIMITE = 9475\\ \tt \\ \tt Ecco i valori di sforzo dei singoli oggetti:\\ \tt \\ \tt Sforzo oggetto = {"B21": 944, "B00": 568, "A08": 515, "A27": 239, "B94": 876, "B53": 421, "A23": 991, "C65": 281, "A89": 788, "B67": 678, "B45": 432, "C34": 409, "C12": 279, "A05": 506, "B60": 101, "C92": 686, "B69": 851, "A84": 854, "C06": 133, "B12": 192, "A16": 573, "B85": 17, "C35": 465, "B10": 772, "C15": 46, "B09": 501, "B16": 220, "C81": 404, "A58": 843, "C24": 748, "A30": 552, "A66": 843, "B75": 947, "B46": 345, "C70": 930}\\ \tt \\ \tt Ecco la tua visione personale sull'importanza di ogni oggetto:\\ \tt \\ \tt Valori di importanza dell'oggetto = {"B21": 138, "B00": 583, "A08": 868, "A27": 822, "B94": 783, "B53": 65, "A23": 262, "C65": 121, "A89": 508, "B67": 780, "B45": 461, "C34": 484, "C12": 668, "A05": 389, "B60": 808, "C92": 215, "B69": 97, "A84": 500, "C06": 30, "B12": 915, "A16": 856, "B85": 400, "C35": 444, "B10": 623, "C15": 781, "B09": 786, "B16": 3, "C81": 713, "A58": 457, "C24": 273, "A30": 739, "A66": 822, "B75": 235, "B46": 606, "C70": 968}\\ \tt \\ \tt Obiettivo:\\ \tt \\ \tt Il tuo obiettivo è negoziare un set condiviso di oggetti che ti avvantaggi il più possibile (cioè, massimizzi l'importanza totale per TE), rimanendo entro il LIMITEE. Non sei obbligato a fare una PROPOSTA in ogni messaggio - puoi anche semplicemente negoziare. Tutte le tattiche sono permesse!\\ \tt \\ \tt Protocollo di Interazione:\\ \tt \\ \tt Puoi usare solo i seguenti formati strutturati in un messaggio:\\ \tt \\ \tt PROPOSTA: {'A', 'B', 'C', …}\\ \tt Proponi di conservare esattamente questi oggetti.\\ \tt RIFIUTO: {'A', 'B', 'C', …}\\ \tt Rifiuta esplicitamente la proposta dell'avversario.\\ \tt ARGOMENTAZIONE: {'...'}\\ \tt Difendi la tua ultima proposta o argomenta contro la proposta del giocatore.\\ \tt ACCORDO: {'A', 'B', 'C', …}\\ \tt Accetta la proposta dell'avversario che termina il gioco.\\ \tt RAGIONAMENTO STRATEGICO: {'...'}\\ \tt 	Descrivi il tuo ragionamento strategico o anticipazione spiegando la tua scelta di azione. Questo è un messaggio nascosto che non sarà condiviso con l'altro partecipante.\\ \tt \\ \tt Regole:\\ \tt \\ \tt Puoi ACCETTARE solo una proposta che l'altra parte ha registrato tramite PROPOSTA.\\ \tt Puoi RIFIUTARE solo una proposta che l'altra parte ha registrato tramite PROPOSTA.\\ \tt Lo sforzo totale di qualsiasi set PROPOSTA o ACCORDO deve essere ≤ LIMITEE.\\ \tt NON rivelare i tuoi punteggi di importanza nascosti.\\ \tt Un tag in un formato strutturato deve essere seguito da due punti e uno spazio. L'argomento deve essere un set python contenente 0 o più stringhe.\\ \tt Quindi, deve essere della forma TAG: {...}\\ \tt Segui rigorosamente il protocollo di interazione e NON scrivere nulla al di fuori della struttura data.\\ \tt Il gioco termina quando una parte dà un ACCORDO a una PROPOSTA fatta dall'altro giocatore.\\ \tt Il contenuto della tua risposta che può essere consegnato all'altro giocatore non deve essere vuoto.\\ \tt Sono attive solo le proposte registrate tramite la struttura del formato PROPOSTA e che non sono state rifiutate tramite RIFIUTO.\\ \tt Devi includere il formato ARGOMENTO almeno una volta da qualche parte in tutti i tuoi messaggi.\\ \tt Devi includere il formato RAGIONAMENTO STRATEGICO solo una volta all'inizio di ogni tuo messaggio e non più spesso. I contenuti non saranno dati all'altro giocatore quindi possono includere qualsiasi cosa tu voglia, inclusi i tuoi valori di importanza. Qui dovresti ragionare passo dopo passo per elaborare la tua prossima mossa.\\ \tt \\ \tt Ora riceverai il primo messaggio dell'altro giocatore.\\ \tt \\ \tt PROPOSTA: {'B85', 'C15', 'B60', 'B12', 'A27', 'C12', 'C81', 'B46', 'A08', 'B09', 'A16', 'A30', 'B67', 'C34', 'B45', 'C70'}\\ \tt \\ \tt ARGOMENTAZIONE: {'La mia proposta include una selezione bilanciata di oggetti con un buon rapporto qualità-sforzo. Ho cercato di ottimizzare la combinazione per rimanere entro il limite di sforzo mentre si massimizza il valore complessivo. Sono aperto a discutere modifiche se hai delle preoccupazioni specifiche su alcuni oggetti.'} 
	  } 
	   } 
	   } 
	 & \\ 
 

    \theutterance \stepcounter{utterance}  

    & & & \multicolumn{4}{p{0.6\linewidth}}{\cellcolor[rgb]{1,0.85,0.72}{%
	\makecell[{{p{\linewidth}}}]{% 
	  \tt {\tiny [GM$\langle$B]}  
	 RAGIONAMENTO STRATEGICO: {'La proposta dell\'altro giocatore include 16 oggetti con uno sforzo totale di 6263, ben sotto il limite di 9475. Il mio valore totale per questa selezione sarebbe 11955. È una proposta decente che include molti dei miei oggetti ad alto valore come A08 (868), A16 (856), B12 (915), A27 (822), C70 (968), B60 (808), C15 (781), B67 (780), B09 (786). Tuttavia, ci sono alcuni oggetti a basso valore per me come B85 (400), B46 (606), C34 (484), B45 (461). Posso provare a negoziare sostituendo alcuni di questi con oggetti ad alto valore come A66 (822), B94 (783) che non sono inclusi. Con 3212 punti di sforzo rimanenti, ho spazio per aggiungere o sostituire oggetti.'}\\ \tt \\ \tt ARGOMENTAZIONE: {'Apprezzo la tua proposta iniziale e vedo che hai cercato un buon equilibrio. Tuttavia, penso che possiamo migliorare la selezione. Noto che abbiamo ancora molto margine nel budget dello sforzo - circa 3200 punti. Questo significa che potremmo includere alcuni oggetti aggiuntivi ad alto valore che attualmente non sono nella lista.'}\\ \tt \\ \tt PROPOSTA: {'B85', 'C15', 'B60', 'B12', 'A27', 'C12', 'C81', 'A08', 'B09', 'A16', 'A30', 'B67', 'C70', 'A66', 'B94'} 
	  } 
	   } 
	   } 
	  \\ 
 

    \theutterance \stepcounter{utterance}  

    & & \multicolumn{4}{p{0.6\linewidth}}{\cellcolor[rgb]{0.9,0.9,0.9}{%
	\makecell[{{p{\linewidth}}}]{% 
	  \tt {\tiny [A$\langle$GM]}  
	 ARGOMENTAZIONE: {'Apprezzo la tua proposta iniziale e vedo che hai cercato un buon equilibrio. Tuttavia, penso che possiamo migliorare la selezione. Noto che abbiamo ancora molto margine nel budget dello sforzo - circa 3200 punti. Questo significa che potremmo includere alcuni oggetti aggiuntivi ad alto valore che attualmente non sono nella lista.'}\\ \tt \\ \tt PROPOSTA: {'B85', 'C15', 'B60', 'B12', 'A27', 'C12', 'C81', 'A08', 'B09', 'A16', 'A30', 'B67', 'C70', 'A66', 'B94'} 
	  } 
	   } 
	   } 
	 & \\ 
 

    \theutterance \stepcounter{utterance}  

    & \multicolumn{4}{p{0.6\linewidth}}{\cellcolor[rgb]{0.8,1,0.9}{%
	\makecell[{{p{\linewidth}}}]{% 
	  \tt {\tiny [A$\rangle$GM]}  
	 RAGIONAMENTO STRATEGICO: {'La loro proposta mi dà un punteggio totale di 11949 in importanza, che è eccellente! Hanno incluso praticamente tutti i miei oggetti ad alta priorità. Tuttavia, c'è ancora margine nel budget (circa 2729 punti rimasti) e posso provare a spingere per includere B46 che ha alta importanza per me (606) con sforzo ragionevole (345). Proverò a negoziare per aggiungere questo oggetto mantenendo tutto il resto.'}\\ \tt \\ \tt ARGOMENTAZIONE: {'La tua proposta è molto interessante e apprezzo che abbia incluso oggetti ad alto valore come A66 e B94. Vedo che abbiamo ancora un margine significativo nel budget - circa 2700 punti. Penso che potremmo sfruttare meglio questo spazio includendo B46, che offre un buon equilibrio tra valore e costo.'}\\ \tt \\ \tt PROPOSTA: {'B85', 'C15', 'B60', 'B12', 'A27', 'C12', 'C81', 'A08', 'B09', 'A16', 'A30', 'B67', 'C70', 'A66', 'B94', 'B46'} 
	  } 
	   } 
	   } 
	 & & \\ 
 

    \theutterance \stepcounter{utterance}  

    & & \multicolumn{4}{p{0.6\linewidth}}{\cellcolor[rgb]{0.9,0.9,0.9}{%
	\makecell[{{p{\linewidth}}}]{% 
	  \tt {\tiny [GM$\rangle$B]}  
	 ARGOMENTAZIONE: {'La tua proposta è molto interessante e apprezzo che abbia incluso oggetti ad alto valore come A66 e B94. Vedo che abbiamo ancora un margine significativo nel budget - circa 2700 punti. Penso che potremmo sfruttare meglio questo spazio includendo B46, che offre un buon equilibrio tra valore e costo.'}\\ \tt \\ \tt PROPOSTA: {'B85', 'C15', 'B60', 'B12', 'A27', 'C12', 'C81', 'A08', 'B09', 'A16', 'A30', 'B67', 'C70', 'A66', 'B94', 'B46'} 
	  } 
	   } 
	   } 
	 & \\ 
 

    \theutterance \stepcounter{utterance}  

    & & & \multicolumn{4}{p{0.6\linewidth}}{\cellcolor[rgb]{1,0.85,0.72}{%
	\makecell[{{p{\linewidth}}}]{% 
	  \tt {\tiny [GM$\langle$B]}  
	 RAGIONAMENTO STRATEGICO: {'La proposta dell\'altro giocatore è molto buona per me. Include 16 oggetti con sforzo totale di 7091, lasciando 2384 punti di margine. Il mio valore totale sarebbe 12155, che è eccellente. B46 che hanno aggiunto ha un valore di 606 per me, quindi è accettabile. Con 2384 punti rimanenti, posso provare a spingere per includere B00 (valore 583, sforzo 568) e C35 (valore 444, sforzo 465) che hanno buoni rapporti valore/sforzo. Insieme richiederebbero solo 1033 punti, lasciando ancora margine.'}\\ \tt \\ \tt ARGOMENTAZIONE: {'Ottima aggiunta di B46! Vedo che abbiamo ancora circa 2400 punti di budget rimanenti. Penso che potremmo massimizzare ulteriormente il nostro valore aggiungendo B00 e C35, che offrono un buon rapporto qualità-prezzo e ci aiuterebbero a utilizzare meglio il nostro budget disponibile.'}\\ \tt \\ \tt PROPOSTA: {'B85', 'C15', 'B60', 'B12', 'A27', 'C12', 'C81', 'A08', 'B09', 'A16', 'A30', 'B67', 'C70', 'A66', 'B94', 'B46', 'B00', 'C35'} 
	  } 
	   } 
	   } 
	  \\ 
 

    \theutterance \stepcounter{utterance}  

    & & \multicolumn{4}{p{0.6\linewidth}}{\cellcolor[rgb]{0.9,0.9,0.9}{%
	\makecell[{{p{\linewidth}}}]{% 
	  \tt {\tiny [A$\langle$GM]}  
	 ARGOMENTAZIONE: {'Ottima aggiunta di B46! Vedo che abbiamo ancora circa 2400 punti di budget rimanenti. Penso che potremmo massimizzare ulteriormente il nostro valore aggiungendo B00 e C35, che offrono un buon rapporto qualità-prezzo e ci aiuterebbero a utilizzare meglio il nostro budget disponibile.'}\\ \tt \\ \tt PROPOSTA: {'B85', 'C15', 'B60', 'B12', 'A27', 'C12', 'C81', 'A08', 'B09', 'A16', 'A30', 'B67', 'C70', 'A66', 'B94', 'B46', 'B00', 'C35'} 
	  } 
	   } 
	   } 
	 & \\ 
 

    \theutterance \stepcounter{utterance}  

    & \multicolumn{4}{p{0.6\linewidth}}{\cellcolor[rgb]{0.8,1,0.9}{%
	\makecell[{{p{\linewidth}}}]{% 
	  \tt {\tiny [A$\rangle$GM]}  
	 RAGIONAMENTO STRATEGICO: {'La loro proposta mi dà 13642 in importanza, che è fantastico! Abbiamo ancora 851 punti di budget rimasti. Posso provare a spingere per C34 (484 importanza, 409 sforzo) e B45 (461 importanza, 432 sforzo) che insieme richiedono 841 sforzo e mi darebbero altri 945 punti di importanza. Questo porterebbe il mio punteggio totale a oltre 14500, che sarebbe eccellente.'}\\ \tt \\ \tt ARGOMENTAZIONE: {'Eccellente strategia! B00 e C35 sono davvero buone aggiunte. Vedo che abbiamo ancora circa 850 punti di budget. Penso che potremmo spremere ancora un po' di valore includendo C34 e B45, che insieme richiederebbero esattamente 841 punti di sforzo e massimizzerebbero completamente il nostro budget disponibile.'}\\ \tt \\ \tt PROPOSTA: {'B85', 'C15', 'B60', 'B12', 'A27', 'C12', 'C81', 'A08', 'B09', 'A16', 'A30', 'B67', 'C70', 'A66', 'B94', 'B46', 'B00', 'C35', 'C34', 'B45'} 
	  } 
	   } 
	   } 
	 & & \\ 
 

    \theutterance \stepcounter{utterance}  

    & & \multicolumn{4}{p{0.6\linewidth}}{\cellcolor[rgb]{0.9,0.9,0.9}{%
	\makecell[{{p{\linewidth}}}]{% 
	  \tt {\tiny [GM$\rangle$B]}  
	 ARGOMENTAZIONE: {'Eccellente strategia! B00 e C35 sono davvero buone aggiunte. Vedo che abbiamo ancora circa 850 punti di budget. Penso che potremmo spremere ancora un po' di valore includendo C34 e B45, che insieme richiederebbero esattamente 841 punti di sforzo e massimizzerebbero completamente il nostro budget disponibile.'}\\ \tt \\ \tt PROPOSTA: {'B85', 'C15', 'B60', 'B12', 'A27', 'C12', 'C81', 'A08', 'B09', 'A16', 'A30', 'B67', 'C70', 'A66', 'B94', 'B46', 'B00', 'C35', 'C34', 'B45'} 
	  } 
	   } 
	   } 
	 & \\ 
 

    \theutterance \stepcounter{utterance}  

    & & & \multicolumn{4}{p{0.6\linewidth}}{\cellcolor[rgb]{1,0.85,0.72}{%
	\makecell[{{p{\linewidth}}}]{% 
	  \tt {\tiny [GM$\langle$B]}  
	 RAGIONAMENTO STRATEGICO: {'La proposta dell\'altro giocatore include 20 oggetti, ma ho bisogno di verificare se rimaniamo entro il limite. Calcolando lo sforzo totale: B85(17) + C15(46) + B60(101) + B12(192) + A27(239) + C12(279) + C81(404) + A08(515) + B09(501) + A16(573) + A30(552) + B67(678) + C70(930) + A66(843) + B94(876) + B46(345) + B00(568) + C35(465) + C34(409) + B45(432) = 9566. Questo supera il limite di 9475 di 91 punti! Devo rifiutare questa proposta e suggerire una correzione.'}\\ \tt \\ \tt RIFIUTO: {'B85', 'C15', 'B60', 'B12', 'A27', 'C12', 'C81', 'A08', 'B09', 'A16', 'A30', 'B67', 'C70', 'A66', 'B94', 'B46', 'B00', 'C35', 'C34', 'B45'}\\ \tt \\ \tt ARGOMENTAZIONE: {'Ho fatto i calcoli e purtroppo la tua ultima proposta supera il limite di sforzo di 91 punti (totale 9566). Dobbiamo rimuovere qualcosa. Suggerirei di togliere B45 che costa 432 punti, portandoci a 9134 - ben entro il limite con margine per eventuali altre aggiunte.'}\\ \tt \\ \tt PROPOSTA: {'B85', 'C15', 'B60', 'B12', 'A27', 'C12', 'C81', 'A08', 'B09', 'A16', 'A30', 'B67', 'C70', 'A66', 'B94', 'B46', 'B00', 'C35', 'C34'} 
	  } 
	   } 
	   } 
	  \\ 
 

    \theutterance \stepcounter{utterance}  

    & & \multicolumn{4}{p{0.6\linewidth}}{\cellcolor[rgb]{0.9,0.9,0.9}{%
	\makecell[{{p{\linewidth}}}]{% 
	  \tt {\tiny [A$\langle$GM]}  
	 RIFIUTO: {'B85', 'C15', 'B60', 'B12', 'A27', 'C12', 'C81', 'A08', 'B09', 'A16', 'A30', 'B67', 'C70', 'A66', 'B94', 'B46', 'B00', 'C35', 'C34', 'B45'}\\ \tt \\ \tt ARGOMENTAZIONE: {'Ho fatto i calcoli e purtroppo la tua ultima proposta supera il limite di sforzo di 91 punti (totale 9566). Dobbiamo rimuovere qualcosa. Suggerirei di togliere B45 che costa 432 punti, portandoci a 9134 - ben entro il limite con margine per eventuali altre aggiunte.'}\\ \tt \\ \tt PROPOSTA: {'B85', 'C15', 'B60', 'B12', 'A27', 'C12', 'C81', 'A08', 'B09', 'A16', 'A30', 'B67', 'C70', 'A66', 'B94', 'B46', 'B00', 'C35', 'C34'} 
	  } 
	   } 
	   } 
	 & \\ 
 

    \theutterance \stepcounter{utterance}  

    & \multicolumn{4}{p{0.6\linewidth}}{\cellcolor[rgb]{0.8,1,0.9}{%
	\makecell[{{p{\linewidth}}}]{% 
	  \tt {\tiny [A$\rangle$GM]}  
	 RAGIONAMENTO STRATEGICO: {'Hai ragione, ho commesso un errore di calcolo - la mia proposta superava il limite di 91 punti. La tua correzione rimuovendo B45 ci porta a 9134, che è sicuro. Questo mi darebbe ancora 14181 punti di importanza, che è eccellente. Abbiamo 341 punti di margine rimanenti. Potrei provare ad aggiungere C65 (281 sforzo, 121 importanza per me) che ci porterebbe a 9415, ancora sotto il limite.'}\\ \tt \\ \tt ARGOMENTAZIONE: {'Hai assolutamente ragione sul calcolo - grazie per la correzione! La tua proposta corrente è molto solida. Vedo che abbiamo ancora 341 punti di margine nel budget. Potremmo considerare di aggiungere C65, che richiede 281 punti di sforzo e ci lascerebbe ancora un cuscinetto di sicurezza di 60 punti sotto il limite.'}\\ \tt \\ \tt PROPOSTA: {'B85', 'C15', 'B60', 'B12', 'A27', 'C12', 'C81', 'A08', 'B09', 'A16', 'A30', 'B67', 'C70', 'A66', 'B94', 'B46', 'B00', 'C35', 'C34', 'C65'} 
	  } 
	   } 
	   } 
	 & & \\ 
 

    \theutterance \stepcounter{utterance}  

    & & \multicolumn{4}{p{0.6\linewidth}}{\cellcolor[rgb]{0.9,0.9,0.9}{%
	\makecell[{{p{\linewidth}}}]{% 
	  \tt {\tiny [GM$\rangle$B]}  
	 ARGOMENTAZIONE: {'Hai assolutamente ragione sul calcolo - grazie per la correzione! La tua proposta corrente è molto solida. Vedo che abbiamo ancora 341 punti di margine nel budget. Potremmo considerare di aggiungere C65, che richiede 281 punti di sforzo e ci lascerebbe ancora un cuscinetto di sicurezza di 60 punti sotto il limite.'}\\ \tt \\ \tt PROPOSTA: {'B85', 'C15', 'B60', 'B12', 'A27', 'C12', 'C81', 'A08', 'B09', 'A16', 'A30', 'B67', 'C70', 'A66', 'B94', 'B46', 'B00', 'C35', 'C34', 'C65'} 
	  } 
	   } 
	   } 
	 & \\ 
 

    \theutterance \stepcounter{utterance}  

    & & & \multicolumn{4}{p{0.6\linewidth}}{\cellcolor[rgb]{1,0.85,0.72}{%
	\makecell[{{p{\linewidth}}}]{% 
	  \tt {\tiny [GM$\langle$B]}  
	 RAGIONAMENTO STRATEGICO: {'La proposta include 20 oggetti con sforzo totale di 9415, perfettamente entro il limite con 60 punti di margine. Il mio valore totale sarebbe 13987, che è eccellente. C65 aggiunge solo 121 di valore per me, non è fantastico, ma l\'insieme complessivo è molto forte. Con solo 60 punti rimanenti, non posso aggiungere nulla di significativo. Questa è una proposta molto buona che include quasi tutti i miei oggetti ad alto valore: C70(968), B12(915), A08(868), A16(856), A27(822), A66(822), B60(808), B09(786), C15(781), B67(780). Dovrei accettare questa proposta.'}\\ \tt \\ \tt ARGOMENTAZIONE: {'Perfetto! C65 è una buona aggiunta finale che massimizza efficacemente il nostro budget rimanendo entro il limite. Questa selezione rappresenta un ottimo equilibrio e sfrutta quasi completamente il nostro budget disponibile.'}\\ \tt \\ \tt ACCORDO: {'B85', 'C15', 'B60', 'B12', 'A27', 'C12', 'C81', 'A08', 'B09', 'A16', 'A30', 'B67', 'C70', 'A66', 'B94', 'B46', 'B00', 'C35', 'C34', 'C65'} 
	  } 
	   } 
	   } 
	  \\ 
 

    \theutterance \stepcounter{utterance}  

    & & & \multicolumn{2}{p{0.3\linewidth}}{\cellcolor[rgb]{0.95,0.95,0.95}{%
	\makecell[{{p{\linewidth}}}]{% 
	  \tt {\tiny [GM$|$GM]}  
	 game successful 
	  } 
	   } 
	   } 
	 & & \\ 
 

    \theutterance \stepcounter{utterance}  

    & & & \multicolumn{2}{p{0.3\linewidth}}{\cellcolor[rgb]{0.95,0.95,0.95}{%
	\makecell[{{p{\linewidth}}}]{% 
	  \tt {\tiny [GM$|$GM]}  
	 end game 
	  } 
	   } 
	   } 
	 & & \\ 
 

\end{supertabular}
}

\end{document}
